\documentclass{dlproblemset}
\usepackage{tikz}
\tikzset{every picture/.style={line width=0.75pt}}

\title{Redes Convolucionais}
\subtitle{Lista de Exercícios}

\begin{document}

\begin{enumerate}[I)]

    \item A rede convolucional LeNET-5 ilustrada na Figura \ref{fig:lenet5} foi desenvolvida por Yann LeCun e publicada em 1989. Sobre a LeNET-5, responda as perguntas abaixo.

    \begin{figure}[H]
        \centering
        \includegraphics[width=0.9\linewidth]{problems/02-cnn/images/LeNET5.jpeg}
        \caption{LeNET-5}
        \label{fig:lenet5}
    \end{figure}
\begin{enumerate}[a)]
    \item Qual o motivo do número 5 no nome dessa rede?
    \item Quais os tamanhos das convoluções usadas?
    \item Qual o tamanho do campo receptivo de uma unidade dos mapas de ativação 14x14 gerados após a camada de \textit{average pooling} (\textit{subsampling})?
    %\item Qual o tamanho do campo receptivo de uma unidade dos mapas de ativação 5x5 gerados após a camada de \textit{average pooling} (\textit{subsampling})?
    \item Quantos parâmetros treináveis tem essa rede?
    \item A maior parte dos parâmetros treináveis está localizada nas camadas convolucionais ou nas totalmente conectadas?
\end{enumerate}

    \item O que é o \textit{stride} em uma camada convolucional?
    \item O que é o \textit{dilation} em uma camada convolucional?
    \item O que é uma convolução transposta?
    \item Por que dizemos que camadas de agrupamento do tipo \textit{Max Pooling} e \textit{Average Pooling} fazem \textit{subsampling} dos mapas de ativação?
    \item Por que usualmente omitimos o número de canais em uma convolução 2D?
    \item Considerando uma convolução 2D 3x3 que conecta dois mapas de ativação com tamanho 224x224 com $C_1$ e $C_2$ canais respectivamente. Como você poderia inicializar os pesos dessa camada usando a inicialização He?

    \item Dada a matriz de input abaixo, e uma rede com 2 camadas convolucionais, calcule a saída da rede com as configurações a seguir. Para a primeira camada convolucional: Padding = 0, Stride = 1; Kernel = $\mathbf{\Theta}_1$; Ativação linear na saída. Para a segunda camada convolucional: Padding = 1; Stride = 1; Kernel = $\mathbf{\Theta}_2$; Ativação ReLU na saída. Considere \textit{dillation} = 0.

    $$\mathbf{X} = \begin{bmatrix}
    7 & 8 & 9 \\
    6 & 5 & 4 \\
    1 & 2 & 3
    \end{bmatrix}
    \qquad
    \mathbf{\Theta}_1 = \begin{bmatrix}
    1 & 0 & -1 \\
    1 & 0 & -1 \\
    1 & 0 & -1
    \end{bmatrix}
    \qquad
    \mathbf{\Theta}_2 = \begin{bmatrix}
    1 & 1 & 1 \\
    0 & 0 & 0 \\
    -1 & -1 & -1
    \end{bmatrix}$$


    \item Dadas as alternativas abaixo, qual das alternativas representa o Kernel mais indicado para identificação de bordas horizontais em uma imagem?
\begin{enumerate}[a)]
    \item $$\begin{bmatrix}
    -1 & -1 & -1 \\
    -1 & 8 & -1 \\
    -1 & -1 & -1
    \end{bmatrix}$$
    \item $$\begin{bmatrix}
    -1 & 0 & 1 \\
    -2 & 0 & 2 \\
    -1 & 0 & 1
    \end{bmatrix}$$
    \item $$\begin{bmatrix}
    0 & -1 & 0 \\
    -2 & 6 & -2 \\
    0 & -1 & 0
    \end{bmatrix}$$
    \item $$\begin{bmatrix}
    -1 & -2 & -1 \\
    0 & 0 & 0 \\
    1 & 2 & 1
    \end{bmatrix}$$
\end{enumerate}


    \item Dada uma imagem de entrada de $32\times32$ pixels, e duas camadas convolucionais, qual o campo receptivo na imagem original que é representado por um único pixel na saída final? Considere: Kernel para as duas camadas ($\mathbf{\Theta}_1 = \mathbf{\Theta}_2$) tem tamanho $5\times5$; Stride = 1; Dillation = 2.

\end{enumerate}

\newpage
\subsection*{Gabarito}
\color{blue}

\begin{enumerate}[I)]
\item
\begin{enumerate}[a.]
    \item O número 5 indica o número de camadas com parâmetros treináveis. São 2 convoluções e 3 totalmente conectadas.

    \item Na primeira camada convolucional foram convoluções 5x5 com 6 canais de saída; Na segunda foram convoluções 5x5 com 16 canais de saída.

    \item Tamanho 6x6. Uma convolução 5x5 tem campo receptivo de tamanho 5x5, se depois dele fizermos um \textit{average pooling} 2x2 teremos \textit{receptive field} de 6x6.

    %\item Tamanho 16x16. Para responder a essa questão você deve perceber que ao aplicar uma conv após fazer pooling 2x2 com stride 2 o campo perceptivo dobra de tamanho.

    \item Na primeira camada: $5x5x1x6 + 6 = 156$; Na segunda camada: $5x5x6x16 + 16 = 2416$; Na terceira camada $400*120 + 120=48120$; Na quarta camada $120*84 + 84=10164$; Na quinta camada $84*10 + 10=850$. Total 61706.

    \item Nas camadas totalmente conectadas.
\end{enumerate}
\item É o quanto o filtro translada entre cada uma das aplicações. Por padrão em convoluções usamos 1, o que significa que o kernel é transladado de uma posição após cada aplicação do filtro.

\item É um parâmetro usado para aumentar o campo receptivo da convolução. Na prática tornamos o kernel esparso completando com zeros.

\item É uma convolução usada para fazer \textit{upsampling} na imagem. É uma forma de aprender a função de \textit{upsampling}.

\item Pois normalmente usamos essas camadas para diminuir a resolução espacial dos mapas de ativação.

\item Pois usualmente o número de canais da convolução é igual ao número de canais dos mapas de ativação.

\item A inicialização He leva em conta o número de parâmetros que contribuem para cada ativação. Como nesse caso estamos falando de $C_1x3x3$ pesos, podemos inicializar com uma normal com média zero e variância $\frac{2}{9C_1}$

\item Resolução Input 3x3 % resolução questão 8 VIII
    Matrizes Iniciais
    $$
    \mathbf{X} = \begin{bmatrix}
    7 & 8 & 9 \\
    6 & 5 & 4 \\
    1 & 2 & 3
    \end{bmatrix}
    \qquad
    \mathbf{\Theta}_1 = \begin{bmatrix}
    1 & 0 & -1 \\
    1 & 0 & -1 \\
    1 & 0 & -1
    \end{bmatrix}
    \qquad
    \mathbf{\Theta}_2 = \begin{bmatrix}
    1 & 1 & 1 \\
    0 & 0 & 0 \\
    -1 & -1 & -1
    \end{bmatrix}
    $$

    Camada 1 (Conv1)
    Input $3\times3$, Kernel $3\times3$, $p=0$, $s=1 \implies$ Saída $1\times1$.
    $$
    \mathbf{C}_1 = \sum(\mathbf{X} \cdot \mathbf{\Theta}_1) + B_1 \quad (B_1 = 0)
    $$
    $$
    \mathbf{C}_1 = (7\cdot1 + 8\cdot0 + 9\cdot-1) + (6\cdot1 + 5\cdot0 + 4\cdot-1) + (1\cdot1 + 2\cdot0 + 3\cdot-1)
    $$
    $$
    \mathbf{C}_1 = (-2) + (2) + (-2) = -2
    $$
    $$
    \mathbf{C}_1 = \begin{bmatrix} -2 \end{bmatrix}
    $$

    Camada 2 (Conv2)
    Input $1\times1$, Kernel $3\times3$, $p=1$, $s=1 \implies$ Saída $1\times1$.
    Primeiro, aplicamos padding na saída de $\mathbf{C}_1$:
    $$
    \mathbf{C}_{1, padded} = \begin{bmatrix}
    0 & 0 & 0 \\
    0 & -2 & 0 \\
    0 & 0 & 0
    \end{bmatrix}
    $$
    Em seguida, aplicamos $\mathbf{\Theta}_2$, $B_2=1$ e ReLU:
    $$
    \text{Soma Conv} = \sum(\mathbf{C}_{1, padded} \cdot \mathbf{\Theta}_2) = (-2 \cdot 0) = 0
    $$
    $$
    \text{Pré ativação} = \text{Soma Conv} + B_2 = 0 + 1 = 1
    $$
    $$
    Output = \text{ReLU}(1) = 1
    $$

    Resultado Final
    $$
    Output = \begin{bmatrix} 1 \end{bmatrix}
    $$

\item Alternativa D % resolução 9 IX

   \item A primeira camada de convolução com kernel $5\times5$ e dilatação $2$, resulta em um \textit{feature map} com campo receptivo de $9\times9$ sobre a imagem de entrada. A segunda camada aplica a mesma configuração de kernel sobre o \textit{feature map} da primeira camada de convolução. Portanto, o campo receptivo final é o campo de visão da camada anterior ($9\times9$) mais a expansão adicional do seu próprio kernel. Resultando em um campo receptivo de $17\times17$. % questão X

\end{enumerate}
\end{document}
